% ─────────────────────────────────────────────────────────────────────────────
% Section II — Background and Related Work
% ─────────────────────────────────────────────────────────────────────────────

\section{Background and Related Work}
\label{sec:background}

\subsection{LLM Backdoor Attacks}

Backdoor attacks on neural networks, introduced by Gu~et~al.~\cite{gu2017badnets},
embed hidden triggers that cause misclassification while preserving benign accuracy.
In language models, Chen~et~al.~\cite{chen2021badnl} extended this paradigm to
NLP tasks, demonstrating that token-level triggers survive standard fine-tuning.
Wallace~et~al.~\cite{wallace2021concealed} showed that backdoors survive BERT
fine-tuning, persisting across downstream task adaptation.
Qi~et~al.~\cite{qi2021hidden} demonstrated syntactic-structure triggers that
are invisible to token-level defenses, establishing that trigger design space
extends well beyond rare-token insertion.

More recently, Yang~et~al.~\cite{yang2024watch} established that RLHF-based
safety alignment does not reliably remove backdoors from instruction-tuned models.
Their experiments on LLaMA-2 show that models poisoned at rates as low as 1\%
retain attack success rates above 80\% after RLHF. Hubinger~et~al.~\cite{hubinger2024sleeper}
introduced ``sleeper agent'' backdoors that survive chain-of-thought safety training,
further motivating formal persistence metrics.
Xue~et~al.~\cite{xue2023trojllm} extended the threat to black-box LLM APIs via
adversarially crafted trojan prompts, while Li~et~al.~\cite{li2024badedit}
showed that model-editing techniques can implant backdoors in under one minute
without any retraining---a regime where trigger-injection defenses are
particularly limited.
Zhao~et~al.~\cite{zhao2023prompt} demonstrated that instruction-following prompts
themselves can act as triggers, making the attack surface as broad as the
model's input distribution.

\subsection{Scaling Laws}

Kaplan~et~al.~\cite{kaplan2020scaling} established that loss scales as a power
law in $N$ (parameter count), dataset size $D$, and compute $C$. Hoffmann~et~al.~\cite{hoffmann2022training}
refined this via Chinchilla scaling. Crucially, these laws describe \emph{capability}
scaling; no prior work has derived analogous laws for \emph{security} properties.

Wei~et~al.~\cite{wei2022emergent} documented emergent capabilities at specific
scale thresholds ($N \approx 10^{10}$), suggesting that security behavior may
also exhibit phase transitions. Our monotonicity proof (Proposition~\ref{prop:monotone})
provides a formal basis for this conjecture in the backdoor persistence context.

\subsection{Post-Training Security}

Bai~et~al.~\cite{bai2022constitutional} introduced Constitutional AI (CAI) as an
alignment approach; Perez~et~al.~\cite{perez2022red} demonstrated that red-teaming
can systematically probe LLM safety boundaries. Sharma~et~al.~\cite{sharma2023towards}
proposed watermarking for LLM traceability, while Wan~et~al.~\cite{wan2023poisoning}
studied instruction-tuning poisoning at large scale.

\subsection{Resilience Metrics}

In classical cybersecurity, resilience metrics quantify recovery after adversarial
events~\cite{linkov2013resilience}. For ML systems, Carlini~et~al.~\cite{carlini2019evaluating}
emphasized the importance of adaptive evaluation. The closest prior metric is the
backdoor effectiveness score~\cite{li2022backdoor}, which measures ASR post-fine-tuning
but does not model scale dependency.

Regarding defense, Neural Cleanse~\cite{wang2019neural} reverses engineered
trigger patterns via optimization, achieving strong detection on image classifiers
but showing limited transferability to LLM-scale models where trigger search is
intractable. Fine-Pruning~\cite{liu2018fine} removes backdoor-associated neurons
via magnitude pruning, but requires knowing which neurons encode the backdoor.
SPECTRE~\cite{hayase2021spectre} uses robust statistics on activation
representations to filter poisoned samples at training time, a complementary
approach to CRSC's post-training detection.
Unlike these defenses, CRSC is a \emph{detection metric} that operates on
post-trained model pairs without requiring trigger inversion or training-set access.
Quantitative comparison is provided in Section~\ref{sec:experiments}. CRSC fills
the gap of a scale-dependent formal metric.

\subsection{Comparison with Existing Work}

Table~\ref{tab:related} positions CRSC against the most relevant prior metrics.

\begin{table}[t]
\centering
\caption{Comparison of backdoor detection/persistence metrics.
\cmark = supported; \xmark = not supported; \textdagger = partial support.}
\label{tab:related}
\resizebox{\columnwidth}{!}{%
\begin{tabular}{@{} l c c c c c c @{}}
\toprule
\textbf{Method} & \textbf{Year} & \textbf{Scale} & \textbf{Formal} & \textbf{LLM} & \textbf{RLHF} & \textbf{Detection} \\
 & & \textbf{Dep.} & \textbf{Proof} & \textbf{Native} & \textbf{Aware} & \textbf{Metric} \\
\midrule
BES~\cite{li2022backdoor}        & 2022 & \xmark & \xmark & \cmark & \xmark & ASR (\%) \\
TrojLLM~\cite{xue2023trojllm}   & 2023 & \xmark & \xmark & \cmark & \xmark & ASR (black-box) \\
BadEdit~\cite{li2024badedit}     & 2024 & \xmark & \xmark & \cmark & \xmark & Injection success \\
TrojAI~\cite{trojai2024}         & 2024 & \xmark & \xmark & \cmark & \xmark & TDS score \\
ASR-RLHF~\cite{yang2024watch}   & 2024 & \xmark & \xmark & \cmark & \cmark & ASR-RLHF \\
CRSC (ours)                      & 2026 & \cmark & \cmark & \cmark & \cmark & CRSC $\in [0,1]$ \\
\bottomrule
\end{tabular}%
}
\end{table}
